\providecommand{\cF}{\mathcal{F}}
\providecommand{\Barch}{\overline{B}^{\mathrm{ch}}}
\section{The brace dg algebra: operadic structure of the universal bulk}
\label{sec:chiral-Hochschild-brace}
\label{ch:brace}% alias for dangling \ref{ch:...}

The chiral Hochschild cochain complex
$C^\bullet_{\mathrm{ch}}(\cA,\cA)$ is the closed-colour object
$Z^{\mathrm{der}}_{\mathrm{ch}}(\cA)$ in the Vol~II
bulk-boundary pair. Its brace operations are the mixed operations of
the two-coloured dioperad $\SCchtop$: closed inputs act on the open
colour $\cA$, and the directional convention forbids open inputs with
closed output. Thus the brace package is not a Dunn-additivity
identification of $\SCchtop$ with a single-colour $E_3$ operad.
After a Drinfeld associator and the two-coloured cobar homotopy of
Theorem~\ref{thm:chiral-higher-deligne} have been fixed, the pair
\[
\bigl(Z^{\mathrm{der}}_{\mathrm{ch}}(\cA),\cA\bigr)
=
\bigl(C^\bullet_{\mathrm{ch}}(\cA,\cA),\cA\bigr)
\]
corepresents bulk actions on the fixed boundary algebra in the
homotopy category of $W(\SCchtop)$-pairs. The shifted
Gerstenhaber bracket on the brace complex is the deformation
dg Lie algebra of the pair.

\begin{theorem}[Universal brace pair with fixed open colour]%
\label{thm:universal-brace-pair}%
\index{universal bulk!cohomological vs chain-level}%
\index{brace dg algebra!associator}%
Let $\cA$ be an $E_1$-chiral algebra on a smooth complex curve,
lying in the Koszul locus. Fix a Drinfeld associator
$\Phi\in\mathrm{GRT}_1(\Bbbk)$, and assume the heptagon hypotheses
and the chain-level contracting homotopy for the two-coloured
$(\SCchtop)^!$-cobar complex used in
Theorem~\ref{thm:chiral-higher-deligne}. Write
$\mathrm{Pair}_{W(\SCchtop)}(\cA)$ for $W(\SCchtop)$-algebra pairs
whose open colour is identified with $\cA$. Then for every
$(B,\cA)\in\mathrm{Pair}_{W(\SCchtop)}(\cA)$ there is a unique
homotopy class of $W(\SCchtop)$-compatible bulk maps
\[
B\longrightarrow Z^{\mathrm{der}}_{\mathrm{ch}}(\cA)
=C^\bullet_{\mathrm{ch}}(\cA,\cA)
\]
intertwining the bulk-to-boundary action. At the chain level a
representative of this class is not canonical: it depends on the
chosen associator $\Phi$ and on the chosen two-coloured cobar
homotopy, and changing either datum changes the representative by a
gauge equivalence in the corresponding convolution dg Lie algebra.
\end{theorem}

\begin{proof}
The homotopy-category assertion is exactly
Theorem~\ref{thm:chiral-higher-deligne}(2), with the closed colour
identified with the chiral derived centre by
Proposition~\ref{prop:derived-center-E2}. Chain representatives are
obtained by applying the associator-fixed strictification of
Remark~\ref{rem:chd-brace-strict-coordinates} to the mixed
$(\SCchtop)^!$-cobar convolution complex. The gauge statement is the
standard deformation theory of operadic twisting morphisms
\cite[Chs.~6, 10]{LV12}: two choices of strictification represent the
same point of $\pi_0$ precisely when they are connected by a
Maurer--Cartan gauge path in
\[
\Hom\bigl((\SCchtop)^!,\End(Z^{\mathrm{der}}_{\mathrm{ch}}(\cA),\cA)\bigr).
\]
\end{proof}

\begin{remark}[Chain representatives and strict terminality]%
Strict terminality of
$(C^\bullet_{\mathrm{ch}}(\cA,\cA),\cA)$ in the unlocalized
chain category would require a distinguished point in the
associator torsor and a distinguished contracting homotopy of the
two-coloured cobar complex. The theorem supplies the invariant
homotopy class; the strict representative is a coordinate choice.
\end{remark}

\begin{theorem}[Chiral Deligne--Tamarkin comparison by colour and filtration]%
\label{thm:chiral-deligne-tamarkin-comparison}%
\index{Deligne--Tamarkin!chiral}%
\index{Deligne--Tamarkin!chiral!chain-level}%
Let $\cA$ be an $E_1$-chiral algebra and put
$Z^{\mathrm{der}}_{\mathrm{ch}}(\cA)=C^\bullet_{\mathrm{ch}}(\cA,\cA)$.
The following three assertions are the Deligne--Tamarkin content used
by this chapter.
\begin{enumerate}
 \item[\textup{(i)}] The cohomology
 $H^\bullet Z^{\mathrm{der}}_{\mathrm{ch}}(\cA)$ carries the
 Gerstenhaber algebra structure induced by the chiral braces. This is
 the homology $H_\ast(E_2)$-algebra attached to the derived centre by
 Proposition~\ref{prop:derived-center-E2}.
 \item[\textup{(ii)}] After fixing
 $\Phi\in\mathrm{GRT}_1(\Bbbk)$, the formal-disc specialization of
 the chiral brace complex carries the chain-level $E_2$ action
 supplied by Tamarkin formality. Under the comparison
 $C^\bullet_{\mathrm{ch,alg}}(\cA)\simeq
 C^\bullet_{\mathrm{ch,geom}}(\cA)$ of
 Proposition~\ref{thm:chiral-hochschild-models-equivalent}, this is
 the associator-fixed action of
 Theorem~\ref{thm:chd-deligne-tamarkin}.
 \item[\textup{(iii)}] For a genuinely ordered $E_1$-chiral algebra,
 a chain-level $E_2$ action extending the open
 $\Conf_\bullet(\R)$ structure is equivalent to a
 $\Phi$-compatible filtered quasi-isomorphism of
 $\Conf_\bullet(\R)$-operadic cochain objects
 \[
 \mathcal E_\cA:=\End^{\mathrm{ch}}_\cA
 \xrightarrow{\;\simeq\;}H_\ast(\mathcal E_\cA).
 \]
 Its obstruction classes are the successive lifting classes in the
 filtered deformation complex
 \[
 \Def_{\Conf(\R)}(\mathcal E_\cA,H_\ast\mathcal E_\cA)
 =
 \Hom_{\Conf(\R)\text{-}\mathrm{coll}}
 \bigl(\overline{B}\mathcal E_\cA,H_\ast\mathcal E_\cA\bigr).
 \]
 Vanishing of these classes, together with the associator
 compatibility, is the exact chain-level hypothesis.
\end{enumerate}
\end{theorem}

\begin{proof}
Item (i) is the brace-to-Gerstenhaber passage for Hochschild
cochains, applied to the chiral Hochschild model of
Proposition~\ref{prop:derived-center-E2}. Item (ii) is
Theorem~\ref{thm:chd-deligne-tamarkin}: the map
$C_\ast(E_2)\to\mathrm{Br}$ is obtained from Tamarkin's formality
theorem after choosing $\Phi$, and the chiral formal-disc model is
identified with the geometric model by
Proposition~\ref{thm:chiral-hochschild-models-equivalent}. Item (iii)
is the obstruction theory for filtered morphisms of operads, written
in the convolution complex controlling morphisms out of the bar
cooperad \cite[Chs.~6, 10]{LV12}. The open colour remains
$\Conf_\bullet(\R)$ throughout; no Dunn tensor product is applied to
$\cA$.
\end{proof}

\begin{proposition}[$E_1$-chiral chain lift and Theorem~H scope]%
\label{rem:e1-chiral-deligne-chain-level-obstruction}%
\index{Deligne--Tamarkin!chiral!$E_1$ chain-level obstruction}%
\index{Massey products!chiral endomorphism operad}%
Let $\mathcal E_\cA=\End^{\mathrm{ch}}_\cA$ be the chiral
endomorphism operad regarded as a filtered
$\Conf_\bullet(\R)$-operadic cochain object. A chain-level
$E_2$ action on $C^\bullet_{\mathrm{ch}}(\cA,\cA)$ that extends the
open ordered-collision structure is equivalent to a
$\Phi$-compatible filtered quasi-isomorphism
$\mathcal E_\cA\simeq H_\ast(\mathcal E_\cA)$ as in
Theorem~\ref{thm:chiral-deligne-tamarkin-comparison}(iii). The
obstruction to extending a partial lift through the next filtration
stage is a class
\[
\mathfrak o_{N+1}(\mathcal E_\cA,\Phi)
\in
H^1\!\left(
F^{N+1}\Def_{\Conf(\R)}(\mathcal E_\cA,H_\ast\mathcal E_\cA)/
F^{N+2}
\right).
\]
For a binary class $r\in H^0(\mathcal E_\cA)$ with $[r,r]=0$, the
first nonlinear obstruction is the Massey coset
\[
\langle r,r,r\rangle
\in
H^1\!\left(F^3\Def_{\Conf(\R)}(\mathcal E_\cA,H_\ast\mathcal E_\cA)/
F^4\right)\Big/\mathrm{Ind}(r),
\]
where $\mathrm{Ind}(r)$ is generated by left and right multiplication
by $r$ on lower filtration cohomology. Theorem~H concentration is not
a consequence of a generic single-colour $E_3$ action here; strict
$E_3^{\mathrm{top}}$ input enters only after the topologisation
hypotheses of Convention~\ref{conv:chd-e3-chiral-meaning}.
\end{proposition}

\begin{proof}
The equivalence with a filtered quasi-isomorphism is the lifting
criterion for operad maps out of a bar cooperad. The obstruction class
is the usual obstruction to solving the Maurer--Cartan equation one
filtration step further in the convolution dg Lie algebra
\cite[Ch.~6]{LV12}. If $[r,r]=0$, the length-three obstruction is the
triple Massey product obtained from choices of null-homotopy for the
binary bracket; changing the null-homotopy adds precisely
$r\cdot H+H\cdot r$. The final sentence is
Convention~\ref{conv:chd-e3-chiral-meaning}: $\SCchtop$ is
two-coloured until explicit inner-stress topologisation has replaced
the chiral direction by a topological one.
\end{proof}

\begin{proposition}[Yangian $r$-matrix obstruction class]%
\label{prop:massey-rrr-yangian-scope}%
\index{Massey products!triple $\langle r,r,r\rangle$ for Yangian}%
\index{Drinfeld associator!rationality and Massey vanishing}%
Let $\mathfrak{g}$ be a finite-dimensional complex simple Lie
algebra, $\Omega\in(\mathfrak{g}\otimes\mathfrak{g})^{\mathfrak{g}}$
the normalised Casimir (quadratic split of the Killing form), and
$r(u)=\Omega/u$ the rational Yangian classical $r$-matrix controlling
$Y_\hbar(\mathfrak{g})$. Let $\mathcal{E}:=\mathrm{End}^{\mathrm{ch}}_{Y_\hbar(\mathfrak{g})}$
be the chiral endomorphism operad of Remark~\ref{rem:e1-chiral-deligne-chain-level-obstruction}.
Then the following statements hold.
\begin{enumerate}
 \item[\textup{(a)}] The classical Yang--Baxter
 equation $[r_{12},r_{13}]+[r_{12},r_{23}]+[r_{13},r_{23}]=0$ gives
 $[r,r]=0$ in the convolution dg Lie algebra of $\mathcal E$.
 Hence the triple Massey obstruction
 $\langle r,r,r\rangle_\mathcal E$ is defined in the quotient
 \[
 H^3(\mathcal E)\Big/
 \bigl(r\cdot H^\bullet(\mathcal E)+H^\bullet(\mathcal E)\cdot r\bigr)
 \]
 after the standard degree shift from the
 $\Conf_\bullet(\R)$-decorated operadic grading.
 \item[\textup{(b)}] Under the Drinfeld--Kohno specialization
 $\mathfrak t_3\to H^\bullet(\mathcal E)$ sending the generators
 $t_{ij}$ to the Casimir insertions $\Omega_{ij}$, the degree-three
 associator coordinate is represented by the line
 \[
 \theta_3=
 [\Omega_{12},[\Omega_{13},\Omega_{23}]]
 \in
 (\mathfrak g^{\otimes 3})^{\mathfrak g}
 \subset H^3(\mathcal E).
 \]
 \item[\textup{(c)}] A choice of Drinfeld associator $\Phi$ supplies
 the length-three filler in the chain-level Deligne--Tamarkin
 morphism. The obstruction to associator independence is the image of
 the coefficient $\sigma_3(\Phi)$ of $\theta_3$ in the quotient of
 part~\textup{(a)}. CYBE defines the Massey coset; it does not decide
 whether the image of $\theta_3$ is zero in the particular chiral
 endomorphism operad.
 \end{enumerate}
\end{proposition}

\begin{proof}
The CYBE is the Maurer--Cartan equation for the binary class $r$ in
the convolution dg Lie algebra; the construction of the triple Massey
coset and its indeterminacy is the standard one
\cite[Ch.~A.3]{LV12}. Kohno's infinitesimal braid Lie algebra
identifies the three-strand Casimir insertions with the images of
$t_{12},t_{13},t_{23}$ \cite{Kohno1987}. Drinfeld's associator
equations and Tamarkin's formality theorem say precisely that a
chosen associator supplies the higher fillers for the chain map
$C_\ast(E_2)\to\mathrm{Br}$ \cite{Drinfeld90QuasiHopf,Tamarkin2003}.
The displayed quotient is therefore the target in which the
associator coordinate is tested after specialization to the chiral
endomorphism operad.
\end{proof}

\begin{proposition}[Indeterminacy quotient for the depth-three associator coordinate]%
\label{prop:massey-yangian-indeterminacy-nonzero}%
\index{Massey products!indeterminacy analysis for Yangian}%
\index{Drinfeld--Kohno Lie algebra!depth filtration on $\mathfrak{t}_3$}%
In the setting of Proposition~\ref{prop:massey-rrr-yangian-scope},
set
\[
 \mathrm{Ind}
 =
 r\cdot H^\bullet(\mathcal{E})+
 H^\bullet(\mathcal{E})\cdot r
 \subset H^3(\mathcal{E}).
\]
Let
\[
\theta_3=[\Omega_{12},[\Omega_{13},\Omega_{23}]]
\in H^3(\mathcal E)
\]
be the image of the depth-three Drinfeld--Kohno Lie word. For any
chosen associator $\Phi$, the first associator-dependent Yangian
Massey coordinate is the class
\[
\omega_3(\Phi)
=
\sigma_3(\Phi)\,[\theta_3]
\in H^3(\mathcal E)/\mathrm{Ind},
\]
where $\sigma_3(\Phi)$ denotes the coefficient of $\theta_3$ after
substituting $t_{ij}\mapsto\Omega_{ij}$. Thus
$\omega_3(\Phi)=0$ if and only if either $\sigma_3(\Phi)=0$ in the
ground field or $\theta_3\in\mathrm{Ind}$ in the specialized chiral
endomorphism operad. The Drinfeld--Kohno calculation controls
$\theta_3$ before specialization; the Yangian endomorphism operad
decides the final quotient.
\end{proposition}

\begin{proof}
The indeterminacy formula is the Massey indeterminacy for a triple
product with all three entries equal to $r$. A Drinfeld associator is
a group-like series in the completed infinitesimal braid Lie algebra;
projecting its logarithm to the line spanned by the nested Lie word
and then applying $t_{ij}\mapsto\Omega_{ij}$ gives the scalar
$\sigma_3(\Phi)$ multiplying $\theta_3$. Passing to the quotient by
$\mathrm{Ind}$ gives the displayed criterion. No injectivity statement
for $\mathfrak t_3\to H^\bullet(\mathcal E)$ is used.
\end{proof}

\begin{proposition}[Yangian associator coefficient through length three]%
\label{prop:yangian-associator-order-3}%
\index{Drinfeld associator!explicit expansion through $\hbar^3$}%
\index{Yangian!chain-level chiral Deligne--Tamarkin input}%
Let $\Phi$ be a Drinfeld associator over a characteristic-zero field
$\Bbbk$. In the completed free Lie algebra on two generators
$x_0,x_1$ write
\[
\log\Phi(x_0,x_1)
=
\lambda_2(\Phi)[x_0,x_1]
+\lambda_{3,0}(\Phi)[x_0,[x_0,x_1]]
+\lambda_{3,1}(\Phi)[x_1,[x_0,x_1]]
+O(\mathrm{length}\ge4).
\]
After the substitution from the infinitesimal braid Lie algebra to
Yangian Casimir insertions, the length-three coordinate used in
Proposition~\ref{prop:massey-yangian-indeterminacy-nonzero} is the
linear combination of $\lambda_{3,0}(\Phi)$ and
$\lambda_{3,1}(\Phi)$ multiplying
$[\Omega_{12},[\Omega_{13},\Omega_{23}]]$; this coefficient is, by
definition, $\sigma_3(\Phi)$. For the KZ associator with the standard
$1/(2\pi i)$ KZ normalization, the corresponding length-three
coefficients lie in
$\Q\cdot \zeta(3)/(2\pi i)^3$ \cite{Drinfeld90QuasiHopf,Kontsevich1999}.
\end{proposition}

\begin{proof}
The logarithm of a group-like associator is Lie. Its length-two and
length-three projections are therefore linear combinations of the
displayed Lie words. Drinfeld computes these coefficients for the KZ
parallel transport on $M_{0,4}$; Kontsevich identifies the odd
length-three period with the motivic $\zeta(3)$ coordinate in the
formality construction. The substitution to Casimir insertions is the
Drinfeld--Kohno representation of $\mathfrak t_3$.
\end{proof}

\begin{proposition}[Rational associators and the $\mathrm{GRT}_1$ torsor]%
\label{prop:yangian-associator-order-3-explicit}%
\index{Drinfeld associator!order-$\hbar^3$ expansion}%
\index{KZ associator!$\zeta(3)$ coefficient}%
\index{A-infinity!associator-parametrised chain-level Deligne--Tamarkin}%
Drinfeld associators over $\Q$ exist. For every characteristic-zero
field $\Bbbk$ for which the set
$\mathrm{Assoc}(\Bbbk)$ is non-empty, it is a torsor under
$\mathrm{GRT}_1(\Bbbk)$; the KZ associator is the distinguished
complex point obtained from regularised KZ parallel transport and has
multiple-zeta coefficients. Consequently:
\begin{enumerate}
 \item[\textup{(i)}] the chain-level chiral Deligne--Tamarkin
 structure on a Yangian chiral endomorphism operad requires a chosen
 associator $\Phi$ as input;
 \item[\textup{(ii)}] changing $\Phi$ acts through the
 $\mathrm{GRT}_1$ torsor on the strict brace coordinates;
 \item[\textup{(iii)}] the associator-independent datum is the
 homotopy class after forgetting these strict coordinates, not a
 canonical rational associator selected by the rational $r$-matrix
 $r(u)=\Omega/u$.
\end{enumerate}
\end{proposition}

\begin{proof}
Drinfeld constructs associators and proves the torsor statement
\cite{Drinfeld90QuasiHopf}; Furusho proves the compatibility of the
pentagon and hexagon equations used in this formulation
\cite{Furusho2010PentagonImpliesHexagon}. Tamarkin's formality
construction uses an associator to produce a chain map from
$C_\ast(E_2)$ to the brace operad \cite{Tamarkin2003}. The rational
Yangian $r$-matrix supplies the infinitesimal braid representation,
but a point of the associator torsor is extra data.
\end{proof}

\begin{remark}[Primary anchors for the associator obstruction]%
\label{rem:massey-rrr-yangian-verification}%
The local obstruction package uses three external inputs.
\begin{enumerate}
 \item Drinfeld's quasi-Hopf theorem supplies associators, the KZ
 point, and the $\mathrm{GRT}_1$ torsor
 \cite{Drinfeld90QuasiHopf}.
 \item Tamarkin's formality theorem supplies the associator-dependent
 chain map from little disks to braces \cite{Tamarkin2003}.
 \item Kohno's infinitesimal braid Lie algebra supplies the
 three-strand Casimir specialization \cite{Kohno1987}; Brown's
 mixed-Tate theorem identifies the motivic odd-depth generators in
 the period algebra \cite{BrownMixedTate}.
\end{enumerate}
\end{remark}

\begin{remark}[$E_1$ primacy and the brace complex]
\label{rem:brace-e1-primacy}
\index{brace dg algebra!E1 primacy@$E_1$ primacy}
The brace operations are the mixed-sector content of the
$\mathsf{SC}^{\mathrm{ch,top}}$ Koszul dual cooperad
(Volume~I, Proposition~\ref*{V1-prop:mixed-sector-bulk-boundary}):
they encode the bulk-to-boundary module structure
$\mu_{k;m}\colon (Z_{\mathrm{ch}}^{\mathrm{der}})^{\otimes k}
\otimes \cA^{\otimes m} \to \cA$
by which the chiral derived center acts on the boundary algebra.
In the $E_1$ primacy framework
(Volume~I, Theorem~\ref*{V1-thm:e1-primacy}),
the brace complex is the mixed sector of the SC convolution algebra,
mediating between the closed factor
$\mathfrak{g}^{\mathrm{mod}}$ (modular shadows) and the open factor
$\mathfrak{g}^{\mathbb{R}}$ ($r$-matrix, Yangian)
(Volume~I, Corollary~\ref*{V1-cor:convolution-factorization}).
\end{remark}

\subsection{Operadic cochains for colored operads and the brace structure}
\label{subsec:operadic-cochains-brace}
Let $O$ be a colored dg--operad and $A$ an $O$--algebra in a symmetric monoidal dg--category $(\mathcal C,\otimes)$.
Write $O^\vee$ for the bar cooperad of $O$ (cofree conilpotent cooperad on $s^{-1}\overline O$).
We adopt the standard operadic cochain complex
\[
\mathrm{C}^\bullet_O(A) \;:=\; \underline{\mathrm{Hom}}_{\text{cooperad}} \bigl( O^\vee,\, \mathrm{End}_A \bigr),
\]
equipped with the differential induced by the coderivation on $O^\vee$ and the internal differential of $\mathrm{End}_A$.
When $O = \OHT$ is our Swiss–cheese chain operad, we denote
\[
\mathrm{C}^\bullet_{\mathrm{ch,top}}(A)\;:=\;\mathrm{C}^\bullet_{\OHT}(A)
\quad\text{(the \emph{chiral Hochschild cochains})}.
\]

\begin{definition}[Brace operations]
\label{def:brace-ops-brace}
The bar cooperad $O^\vee$ carries a canonical \emph{brace} structure; by precomposition,
$\mathrm{C}^\bullet_O(A)$ acquires multilinear operations
\[
\{-\}\{-,\ldots,-\}\colon \mathrm{C}^\bullet_O(A)\otimes \mathrm{C}^\bullet_O(A)^{\otimes m}\longrightarrow \mathrm{C}^\bullet_O(A),
\]
given by grafting trees in $O^\vee$ (see \cite{GJ94,Vor99}).
The resulting algebraic structure on $\mathrm{C}^\bullet_O(A)$ is a homotopy Gerstenhaber (brace) algebra; its induced \emph{Gerstenhaber bracket} is
\[
[f,g] := f\{g\} - (-1)^{(\deg f-1)(\deg g-1)}\, g\{f\}.
\]
\end{definition}

\begin{remark}[Braces as strict coordinates]
\label{rem:brace-strict-coordinates}
\index{brace algebra!as strict coordinates}
The brace algebra is not the invariant deformation object; it is a
convenient strict presentation of the $E_2$/Gerstenhaber package on
Hochschild cochains. Equivalently, the Gerstenhaber bracket on
$\mathrm{C}^\bullet_O(A)[1]$ is the binary Taylor coefficient of the
filtered $L_\infty$-structure obtained from the corresponding
bar-coderivation complex.
\end{remark}

\begin{remark}[Chiral brace vs topological brace]%
\label{rem:brace-chiral-vs-topological-discipline}%
\index{brace operations!chiral vs topological}%
\index{brace operations!Hochschild-theory discipline}%
The brace operations on
$\mathrm{C}^\bullet_{\mathrm{ch,top}}(A)$ are
\emph{chiral Gerstenhaber brace operations}, distinct from classical
\emph{topological Gerstenhaber brace} operations. The discipline:
\begin{itemize}[nosep]
\item (Three Hochschild theories, one brace each.) The brace acts on
 $\mathrm{ChirHoch}^\bullet$, rather than on
 $\HH^\bullet(A_{\mathrm{mode}})$ (topological Deligne brace) or on
 $\HH^\bullet_{\mathrm{cat}}$ (categorical CY shifted-Poisson brace).
 Each Hochschild theory carries
 its own brace; they coincide only after the comparison maps of
 Theorem~\ref{thm:three-hochschild-unification} (\S\ref{sec:chiral_hochschild}).
\item (Chiral $\ne$ topological brace.) The chiral brace is
 defined via integration over $\FM_k(\C) \times \Conf_m(\R)$
 with logarithmic kernels; the topological brace is the McClure--Smith
 brace on the little 2-disks operad. They agree on
 $E_\infty$-chiral algebras after Tamarkin formality and a Drinfeld
 associator (Theorem~\ref{thm:chiral-deligne-tamarkin-comparison}(ii));
 on genuinely $E_1$-chiral algebras (Yangian, EK quantum
 vertex algebras), strict identification is governed by the
 obstruction quotient of Proposition~\ref{prop:massey-rrr-yangian-scope}.
\item The unqualified phrase ``Hochschild brace'' is ambiguous here;
 the intended brace is always chiral, topological, or categorical.
\end{itemize}
\end{remark}

\begin{remark}[Degree conventions]
We use cohomological grading. For $O=\OHT$ the internal degrees reflect the form‑degrees on $\FM(\C)$ and the $E_1$‑degrees on $\Conf(\R)$; the brace degree shifts match the usual Gerstenhaber sign rule.
The brace operations are the explicit transferred $L_\infty$-brackets
from the operadic cochain complex, realising $\Convinf$ via tree-grafting
by the operadic homotopy-convolution theorem of Volume~I.
\end{remark}

\subsection{\texorpdfstring{Geometric brace model on $\FM(\C)\times\Conf(\R)$}{Geometric brace model on FM(C) x Conf(R)}}
\label{subsec:geometric-brace-expanded}
The brace operations admit a concrete model via configuration spaces.
Let $\Conf_k(\R)$ denote ordered configurations $(t_1<\cdots<t_k)$ of $k$ points in~$\R$ and set
\[
\mathbb{X}_{k,m} \;:=\; \FM_k(\C)\times \Conf_m(\R).
\]
Operations in $\OHT$ at degree $(k,m)$ are represented by chains on $\mathbb{X}_{k,m}$ decorated by edge‑length parameters.
Given $f\in \mathrm{C}^\bullet_{\mathrm{ch,top}}(A)$ and $g_1,\dots,g_m\in \mathrm{C}^\bullet_{\mathrm{ch,top}}(A)$, we define
\[
f\{g_1,\dots,g_m\} := \int_{\mathbb{X}} \;\;\Bigl(\text{pullback of kernels from } \mathbb{X}_{k,m'} \Bigr)\;\circ\;\Bigl(\text{insertions via trees along } \FM\times \Conf\Bigr),
\]
where the integral is over the appropriate fiber product of configuration spaces determined by the brace tree (the precise integral formula with signs is given in Appendix~\ref{app:brace-signs}; the geometric content is the operadic composition of correlation kernels via fiber integration over the tree parameter space).

\begin{proposition}[Well-definedness and functoriality]
\label{prop:geometric-braces-well-defined-brace}
If $A$ is a logarithmic $\SCchtop$-algebra, then the geometric brace operations on $\mathrm{C}^\bullet_{\mathrm{ch,top}}(A)$ are well-defined, strictly multilinear, and satisfy the brace identities. The Gerstenhaber bracket obtained from these braces coincides with the one induced by the cooperad $O^\vee$.
\end{proposition}

\begin{proof}
Integrability follows from Definition~\ref{def:log-SC-algebra} (logarithmic weight forms have controlled singularities): the weight forms have at most logarithmic singularities along the boundary divisors of the FM compactification, and such forms are integrable on compact manifolds with corners. The brace identities follow from Stokes' theorem on FM boundaries with corner cancellation by Theorem~\ref{thm:Arnold_full_proof} (Appendix~\ref{app:FM_Stokes}): the boundary strata reproduce the algebraic brace relations, with the Arnold cancellations at codimension-2 corners ensuring that no spurious terms arise (see Appendix~\ref{app:brace-signs} for the detailed sign verification). Compatibility with the cooperad differential is the standard operadic cochain construction for a $W(\SCchtop)$-algebra.
\end{proof}

\subsection{The \texorpdfstring{$\ast$}{*}–Lie structure and deformation theory}
\label{subsec:star-Lie-brace}
Let $(\mathrm{C}^\bullet_{\mathrm{ch,top}}(A),\{-\}\{-\})$ be the brace algebra constructed above.

\begin{definition}[$\ast$–Lie algebra]
\label{def:star-Lie-brace}
Define the shifted Lie bracket on $\mathrm{C}^\bullet_{\mathrm{ch,top}}(A)[1]$ by the Gerstenhaber bracket $[-,-]$ associated to braces. We call $(\mathrm{C}^\bullet_{\mathrm{ch,top}}(A)[1],[-,-])$ the \emph{chiral $\ast$–Lie algebra} of $A$.
\end{definition}

\begin{theorem}[Deformations and Maurer–Cartan]
\label{thm:MC-deformations-brace}
Let $A$ be a $\OHT$–algebra, and consider formal deformations over $k[\![\hbar]\!]$. Then there is a bijection between:
\begin{enumerate}[label=(\alph*)]
 \item deformation classes of $A$ as a $W(\SCchtop)$–algebra over $k[\![\hbar]\!]$ up to gauge; and
 \item gauge‑equivalence classes of solutions to the Maurer–Cartan equation
 \(
 d\alpha + \tfrac{1}{2}[\alpha,\alpha]=0
 \)
 in $\bigl(\mathrm{C}^\bullet_{\mathrm{ch,top}}(A)[1]\otimes \hbar k[\![\hbar]\!],[-,-]\bigr)$.
\end{enumerate}
\end{theorem}

\begin{proof}
By the Kontsevich--Soibelman deformation theorem \cite[Theorem~3.1]{KS09} (see also \cite[Theorem~10.3.1]{LV12}), deformations are classified by twisting morphisms from the bar cooperad into $\mathrm{End}_A$, and a twisting morphism is equivalent to an MC element in the homotopy convolution object associated to $(O^\vee,\mathrm{End}_A)$. On a chosen bar-cofree model this homotopy object is strictified by the convolution Lie algebra $\mathrm{Hom}(O^\vee,\mathrm{End}_A)$; see \cite{GK94,Val07}. The brace model identifies this strictification with $\mathrm{C}^\bullet_{\mathrm{ch,top}}(A)[1]$.
\end{proof}

\subsection{Bulk observables \texorpdfstring{$\simeq$}{≃} chiral Hochschild cochains}
\label{subsec:bulk-equals-CHC-brace}
Let $\mathsf{Obs}$ be the BV observable prefactorization algebra of the HT theory, and $A=\mathsf{Obs}(D\times I)$ the local complex on a small rectangle.

\begin{theorem}[Bulk--cochains identification under product-formal HT hypotheses]
\label{thm:bulk-CHC-brace}
Let $\mathsf{Obs}$ be an HT prefactorization algebra satisfying the
physics-bridge hypotheses, including a product-formal local
$\SCchtop$ shadow on small rectangles, and put
$A=\mathsf{Obs}(D\times I)$. Then there is a canonical filtered
quasi-isomorphism
\[
\mathrm{C}^\bullet_{\mathrm{ch,top}}(A)\;\simeq\;\mathsf{Obs}_{\mathrm{bulk}}(\mathbb{B}^3)
\]
between chiral Hochschild cochains of the local boundary algebra $A$
and the product-formal local bulk cochains supported in a small
$3$-ball, with filtration by holomorphic weight. On associated
graded, this identification reduces to the BD-chiral Hochschild complex
on the closed color tensored with the $E_1$ Hochschild complex on the
open color. An arbitrary logarithmic $\SCchtop$-algebra enters this
theorem only after a chosen HT prefactorization realization supplies
the product-formal local shadow.
\end{theorem}

\begin{proof}
The argument has three steps: product-formal operad-of-shapes
recognition providing a common local model, filtration analysis via
holomorphic weight, and integrability from the chosen HT
prefactorization realization.

\medskip
\textbf{Step 1: Operad-of-shapes recognition.}

Decompose the small $3$--ball $\mathbb{B}^3 \subset \C \times \R$ into rectangles $D_i \times I_j$ (discs in $\C$ times intervals in $\R$). On each rectangle, $\mathsf{Obs}(D_i \times I_j) \simeq A$ by local constancy (Proposition~\ref{prop:prefact-Obs}). The factorization homology $\int_{\mathbb{B}^3} \mathsf{Obs}$ is the homotopy colimit over all such decompositions; the indexing category is the \emph{operad of shapes} $\mathbb{X}_{k,m} = \FM_k(\C) \times \Conf_m(\R)$, where $k$ holomorphic and $m$ ordered topological insertion points parametrize embeddings of rectangles in $\mathbb{B}^3$.

An element of $\mathsf{Obs}_{\mathrm{bulk}}(\mathbb{B}^3)$ is a compatible family of multilinear operations
\[
A^{\otimes k} \otimes A^{\otimes m} \longrightarrow A,
\]
parametrized by points in $\mathbb{X}_{k,m}$, i.e., by integration of propagator kernels over $\FM_k(\C) \times \Conf_m(\R)$. By the Swiss--cheese recognition (Theorem~\ref{thm:recognition-foundations}), these operations assemble into a $W(\SCchtop)$--algebra, and the totality of such operations is precisely the operadic cochain complex
\[
\mathrm{C}^\bullet_{\mathrm{ch,top}}(A) = \underline{\mathrm{Hom}}_{\text{cooperad}}\bigl((\OHT)^\vee, \mathrm{End}_A\bigr).
\]
This gives the map $\Phi \colon \mathrm{C}^\bullet_{\mathrm{ch,top}}(A) \to \mathsf{Obs}_{\mathrm{bulk}}(\mathbb{B}^3)$: a cochain $f$ of degree $(k,m)$ determines an observable by integrating $f$ against the propagator kernels over $\mathbb{X}_{k,m}$. In the reverse direction, restricting a bulk observable to the configuration-space strata recovers the cochain components.

\medskip
\textbf{Step 2: Holomorphic weight filtration.}

Equip both sides with the \emph{holomorphic weight filtration} $F^p$, graded by total pole order at collision divisors in $\FM_k(\C)$.

On the cochain side, define $F^p \mathrm{C}^\bullet_{\mathrm{ch,top}}(A) \subset \mathrm{C}^\bullet_{\mathrm{ch,top}}(A)$ as the subcomplex of cochains whose kernels, when restricted to $\FM_k(\C)$, have poles of total order at most $p$ along the union of collision divisors $\bigcup_{|S| \geq 2} D_S$.

On the observable side, define $F^p \mathsf{Obs}_{\mathrm{bulk}}(\mathbb{B}^3)$ by the analogous pole-order condition on the Feynman-diagrammatic integrands.

The map $\Phi$ is filtered: it preserves pole order since it is defined by integration over the \emph{same} configuration-space kernels. On the associated graded, the factorization structure splits by the product decomposition $\mathbb{X}_{k,m} = \FM_k(\C) \times \Conf_m(\R)$. At each graded piece:
\begin{enumerate}[label=(\roman*)]
\item The \emph{chiral component} (varying $k$, from $\FM_k(\C)$) contributes the BD--chiral Hochschild complex $\mathrm{C}^\bullet_{\mathrm{BD\text{-}ch}}(A_{\mathsf{ch}})$, where $A_{\mathsf{ch}}$ is the closed-color (holomorphic) part of the $W(\SCchtop)$--algebra. The differential on this piece is the chiral Hochschild differential with $\delta_Q$ from the internal BV-BRST operator and $\delta_{\mathrm{Hoch}}$ from the boundary faces of $\FM_k(\C)$ encoding OPE collisions.
\item The \emph{topological component} (varying $m$, from $\Conf_m(\R)$) contributes the $E_1$ Hochschild complex $\mathrm{C}^\bullet_{E_1}(A_{\mathsf{top}})$, where $A_{\mathsf{top}}$ is the open-color (topological) part. The differential encodes the boundary faces of $\Conf_m(\R)$ (merging adjacent intervals).
\end{enumerate}
The splitting on associated graded is
\[
\gr^F \mathrm{C}^\bullet_{\mathrm{ch,top}}(A) \;\cong\; \mathrm{C}^\bullet_{\mathrm{BD\text{-}ch}}(A_{\mathsf{ch}}) \;\widehat\otimes\; \mathrm{C}^\bullet_{E_1}(A_{\mathsf{top}}),
\]
and the same decomposition holds for $\gr^F \mathsf{Obs}_{\mathrm{bulk}}(\mathbb{B}^3)$ by the factorization product structure on disjoint rectangles. Hence $\gr(\Phi)$ is an isomorphism.

\medskip
\textbf{Step 3: Well-definedness and integrability.}

The integrals defining $\Phi$ converge because Definition~\ref{def:log-SC-algebra} supplies kernels on $\FM_k(\C) \times \Conf_m(\R)$ whose holomorphic part has only logarithmic singularities, and Theorem~\ref{thm:FM-calculus} provides the residue and boundary control needed for the compactified integral. Since $\FM_k(\C)$ is a compact manifold with corners, forms with logarithmic singularities are integrable; the key estimate is
\[
\int_0^\epsilon |\log \varepsilon_S|^N \, \varepsilon_S^{a-1} \, d\varepsilon_S < \infty \quad \text{for } a > 0, \; N \geq 0,
\]
where $\varepsilon_S$ is the radial parameter measuring the size of the cluster $S$ (Definition~\ref{def:FM_coords}). The exponent $a > 0$ is exactly the integrability margin furnished by the logarithmic singularity condition. For physical realizations, the familiar meromorphic propagator and topological decay statements are supplied by Theorem~\ref{thm:physics-bridge}.

The filtration is complete and Hausdorff on both sides (pole order is bounded below and the complexes are exhausted by finite filtration levels). By the spectral sequence comparison theorem, the isomorphism on $\gr$ lifts to a quasi-isomorphism on the filtered complexes:
\[
\Phi \colon \mathrm{C}^\bullet_{\mathrm{ch,top}}(A) \;\xrightarrow{\;\simeq\;}\; \mathsf{Obs}_{\mathrm{bulk}}(\mathbb{B}^3). \qedhere
\]
\end{proof}

\begin{corollary}[Obstruction theory]
\label{cor:obstructions-brace}
For the controlling dg Lie algebra
$\mathfrak g_A=\mathrm{C}^\bullet_{\mathrm{ch,top}}(A)[1]$,
infinitesimal deformations are parameterized by
$H^1(\mathfrak g_A)$ and primary obstructions lie in
$H^2(\mathfrak g_A)$. Equivalently, in the unshifted cochain grading,
infinitesimals lie in
$H^2(\mathrm{C}^\bullet_{\mathrm{ch,top}}(A))$ and primary
obstructions lie in
$H^3(\mathrm{C}^\bullet_{\mathrm{ch,top}}(A))$. Higher obstructions
live in the Massey--product tower controlled by the brace operations.
\end{corollary}

\subsection{Filtration and PVA on associated graded}
\label{subsec:filtration-PVA-brace}
Let $F^\bullet$ be the holomorphic weight filtration on $\mathrm{C}^\bullet_{\mathrm{ch,top}}(A)$. Then
\[
\mathrm{gr}^F\, \mathrm{C}^\bullet_{\mathrm{ch,top}}(A)
\;\cong\;
\mathrm{C}^\bullet_{\mathrm{BD\text{-}ch}}(A_{\mathsf{ch}})\;\widehat\otimes\; \mathrm{C}^\bullet_{E_1}(A_{\mathsf{top}}),
\]
and the induced bracket on $\mathrm{H}^\bullet$ coincides with the $(-1)$‑shifted Poisson vertex bracket of Section~\ref{sec:Ainfty-to-PVA}.
